
Áreas metropolitanas y urbanas de Nuevo León con gran movimiento industrial han presentado niveles de O₃ y PM10 en su atmósfera que, en diversas ocasiones, han llegado a superar los límites permisibles establecidos por las normas oficiales mexicanas correspondientes: la NOM-020-SSA1-2021 para ozono y la NOM-025-SSA1-2014 para partículas PM10 y PM2.5. Estas tendencias de contaminación del aire se han comenzado a documentar desde hace más de 30 años, con distintas herramientas de detección que permiten medir tanto las concentraciones de ozono (mediante fotometría ultravioleta) como las de sus precursores, los óxidos de nitrógeno NO y NO₂ (agrupados como NOx, medidos por quimioluminiscencia), cuya relación fotoquímica es clave para entender la formación del contaminante secundario O₃.

El incremento sostenido del ozono a nivel del suelo se puede atribuir a la expansión de la ciudad, tanto industrial como poblacional, tomando en cuenta el aumento acelerado del parque vehicular que ha habido en los últimos años. El estado meteorológico juega un papel fundamental en la distribución del O₃ a nivel del suelo, un papel que se ha visto adicionalmente modulado por el cambio climático en los últimos años.

Diversos estudios han evidenciado que el Área Metropolitana de Monterrey no es la excepción a este fenómeno: entre 1993 y 2014 se registró un aumento sostenido de las concentraciones de ozono troposférico, en contraste con la tendencia descendente observada en la Ciudad de México durante el mismo periodo, lo cual ha sido atribuido al crecimiento industrial y vehicular de la región sin un control equivalente sobre sus precursores atmosféricos. Este comportamiento se agrava por las condiciones topográficas del valle donde se asienta el AMM, rodeado de sierras que favorecen el estancamiento de masas de aire durante episodios de alta presión atmosférica (anticiclones), inversión térmica y baja velocidad de viento, factores que reducen la dispersión de contaminantes y prolongan los periodos de exposición de la población a concentraciones elevadas. 

